% ============================================================
%  CAPÍTULO 3 — PLANIFICACIÓN
% ============================================================
\chapter{Planificación}
\label{ch:planificacion}

\section{Metodología de trabajo}
\label{sec:metodologia}

El desarrollo del proyecto se ha realizado siguiendo una metodología \textbf{iterativa e
incremental} adaptada a la naturaleza del trabajo: un único desarrollador con reuniones de
seguimiento periódicas con el tutor de empresa y la tutora académica.

Cada iteración tiene una duración aproximada de tres semanas y produce un incremento
funcional y demostrable. Al inicio de cada iteración se identifican las tareas a abordar;
al final se realiza una revisión con el tutor de empresa para validar el incremento y
priorizar la siguiente iteración.

\subsection{Herramientas de gestión}
\label{subsec:herramientas-gestion}

\begin{itemize}
  \item \textbf{Control de versiones:} Git con repositorios en GitHub (repositorio
    privado en la organización del \gls{itc}).
  \item \textbf{Seguimiento de tareas:} tablero Kanban propio con las columnas
    \textit{Backlog}, \textit{En progreso} y \textit{Hecho}.
  \item \textbf{Documentación técnica:} comentarios \textit{docstring} en el código y
    este documento de memoria.
  \item \textbf{Comunicación:} reuniones semanales por videoconferencia con el equipo del
    \gls{itc} y tutorías presenciales o telemáticas con la tutora académica.
\end{itemize}

\section{Fases e iteraciones}
\label{sec:fases}

El desarrollo se ha organizado en cuatro iteraciones principales, descritas en la
tabla~\ref{tab:iteraciones}.

\begin{table}[H]
  \centering
  \caption{Iteraciones del proyecto y sus objetivos principales.}
  \label{tab:iteraciones}
  \begin{tabular}{clp{7cm}}
    \toprule
    \textbf{Iter.} & \textbf{Período (aprox.)} & \textbf{Objetivos} \\
    \midrule
    1 & Enero -- Feb.~2026 &
      Análisis del código existente; tubería SQLAlchemy~→~JSON~Schema;
      endpoints de pantallas dinámicas (\textit{screens}). \\
    2 & Feb. -- Mar.~2026 &
      Sistema de plugins por capas; integración frontend-backend para
      tablas dinámicas y formularios básicos. \\
    3 & Mar. -- Abr.~2026 &
      Generación de formularios con Angular~Formly; endpoint
      \textit{bulk}; cuadrícula de edición masiva (AG~Grid). \\
    4 & Abr. -- Jun.~2026 &
      Editor visual de esquemas (Formly~Editor); pruebas de integración;
      documentación y pulido final. \\
    \bottomrule
  \end{tabular}
\end{table}

\section{Desglose de tareas (WBS)}
\label{sec:wbs}

A continuación se detallan las tareas principales agrupadas por área funcional.

\subsection{Backend — uniback}
\label{subsec:wbs-backend}

\begin{enumerate}[label=B\arabic*.]
  \item Análisis y comprensión del código base existente.
  \item Diseño de la tubería de transformación SQLAlchemy~→~Pydantic~→~JSON~Schema.
  \item Implementación de los anotadores de extensiones \texttt{x-*} (tipo, etiqueta,
    opciones, validadores, \ldots).
  \item Sistema de plugins por capas: interfaz, registro y aplicación.
  \item Endpoint \texttt{GET /screens/<nombre>}: generación dinámica de definiciones de
    pantalla.
  \item Endpoint \texttt{POST /<entidad>/bulk}: operaciones en bloque con modos
    \textit{partial} y \textit{abort}.
  \item Cobertura de pruebas: unitarias para la tubería y de integración para los
    endpoints.
\end{enumerate}

\subsection{Frontend — ngt-gui}
\label{subsec:wbs-frontend}

\begin{enumerate}[label=F\arabic*.]
  \item Integración de Angular~Formly en la librería.
  \item Mapeo JSON~Schema~→~configuración Formly (\texttt{schema-to-formly.ts}).
  \item Componente \texttt{DynamicFormComponent}: formularios generados a partir de
    esquemas.
  \item Componente \texttt{DynamicPageComponent}: enrutamiento y selección de vista
    según \texttt{screen\_type}.
  \item Mapeo JSON~Schema~→~definiciones de columna AG~Grid
    (\texttt{schema-to-aggrid.ts}).
  \item Componente \texttt{BulkEditGridComponent}: cuadrícula de edición masiva tipo
    Excel.
  \item Editor visual de esquemas (\textit{Formly~Editor}): árbol de campos, panel de
    propiedades y vista previa en vivo.
\end{enumerate}

\subsection{Infraestructura y documentación}
\label{subsec:wbs-infra}

\begin{enumerate}[label=I\arabic*.]
  \item Ficheros Docker~Compose para desarrollo y producción.
  \item Configuración de variables de entorno y secretos.
  \item Redacción de la memoria del \gls{tfg}.
  \item Preparación de la presentación oral.
\end{enumerate}

\section{Estimación de recursos}
\label{sec:recursos}

\subsection{Tiempo}
\label{subsec:tiempo}

El \gls{tfg} tiene asignados \textbf{12~ECTS}, equivalentes a aproximadamente 300~horas de
trabajo del estudiante. La distribución estimada se recoge en la
tabla~\ref{tab:horas}.

\begin{table}[H]
  \centering
  \caption{Distribución estimada de horas por área.}
  \label{tab:horas}
  \begin{tabular}{lrr}
    \toprule
    \textbf{Área} & \textbf{Horas estimadas} & \textbf{\%} \\
    \midrule
    Análisis y diseño & 40 & 13 \\
    Implementación backend & 80 & 27 \\
    Implementación frontend & 90 & 30 \\
    Pruebas & 30 & 10 \\
    Documentación (memoria) & 50 & 17 \\
    Presentación & 10 & 3 \\
    \midrule
    \textbf{Total} & \textbf{300} & \textbf{100} \\
    \bottomrule
  \end{tabular}
\end{table}

\subsection{Herramientas de desarrollo}
\label{subsec:herramientas-dev}

El desarrollo se ha realizado íntegramente en un equipo MacBook Pro (Apple~M2,
16~GB~RAM) con el siguiente \textit{software} principal:

\begin{itemize}
  \item Visual Studio Code con las extensiones Angular Language Service, Python,
    Pylance, Docker y LaTeX~Workshop.
  \item Python~3.12 (entorno virtual \textit{venv}).
  \item Node.js~22~LTS con Angular~CLI~19.
  \item PostgreSQL~16 (contenedor Docker para desarrollo local).
  \item Docker Desktop.
\end{itemize}

\section{Diagrama de Gantt}
\label{sec:gantt}

% TODO: insertar figura del diagrama de Gantt
% \begin{figure}[H]
%   \centering
%   \includegraphics[width=\textwidth]{Media/Planificacion/gantt.png}
%   \caption{Diagrama de Gantt del proyecto.}
%   \label{fig:gantt}
% \end{figure}

\todo[inline]{Insertar diagrama de Gantt generado con GanttProject o LaTeX/pgfgantt.}

El diagrama de Gantt refleja la superposición parcial de las fases de análisis,
implementación y documentación, característica de las metodologías iterativas en las que
la documentación se va produciendo en paralelo con el desarrollo.
